\chapter{2025年北京卷}
\section{单选题}

\begin{problem}
集合 \(M = \{x\mid 2x - 1 > 5\}\),\(N=\{1,2,3\}\),则 \(M\cap N=\)(\quad)

\choices
    {\(\{1,2,3\}\)}
    {\(\{2,3\}\)}
    {\(\{3\}\)}
    {\(\emptyset\)}
\end{problem}

\begin{problem}
已知复数 \(z\) 满足 \(\i\cdot z +2 = 2\i\),则 \(|z| =\)(\quad)
\choices
    {\(\sqrt{2}\)}
    {\(2\sqrt{2}\)}
    {4}
    {8}
\end{problem}

\begin{problem}
双曲线 \(x^2-4y^2=4\) 的离心率为(\quad)
\choices
    {\(\dfrac{\sqrt{3}}{2}\)}
    {\(\dfrac{\sqrt{5}}{2}\)}
    {5}
    {\(\sqrt{5}\)}
\end{problem}

\begin{problem}
为得到函数 \(y=9^x\) 的图象,只需把函数 \(y=3^x\) 的图象上的所有点

\choices
    {\(x\) 变成原来的 \(\dfrac12\) 倍,\(y\) 不变}
    {\(x\) 变成原来的 \(2\) 倍,\(y\) 不变}
    {\(y\) 变成原来的 \(\dfrac13\) 倍,\(x\) 不变}
    {\(y\) 变成原来的 \(3\) 倍,\(x\) 不变}
\end{problem}

\begin{problem}
已知 \(\{a_n\}\) 是公差不为 0 的等差数列,\(a_1=-2\),若 \(a_3\),\(a_4\),\(a_6\) 成等比数列,则 \(a_{10}=\)(\quad)
\choices
    {\(-20\)}
    {\(-18\)}
    {16}
    {18}
\end{problem}

\begin{problem}
已知 \(a>0\),\(b>0\),则(\quad)
\choices
    {\(a^2+b^2>2ab\)}
    {\(\dfrac{1}{a}+\dfrac{1}{b}\ge \dfrac{2}{\sqrt{ab}}\)}
    {\(a+b>\sqrt{ab}\)}
    {\(\dfrac{1}{a}+\dfrac{1}{b}\le \dfrac{2}{\sqrt{ab}}\)}
\end{problem}

\begin{problem}
已知函数 \(f(x)\) 的定义域为 \(D\),则``函数 \(f(x)\) 的值域为 \(\R\)''是``对任意 \(M\in\R\),存在 \(x_0\in D\) 使得 \(|f(x_0)|>M\)''的(\quad)
\choices
    {充分不必要条件}
    {必要不充分条件}
    {充分必要条件}
    {既不充分也不必要条件}
\end{problem}

\begin{problem}
设函数 \(f(x)=\sin(\omega x)+\cos(\omega x)\)(\(\omega>0\)),若 \(f(x+\pi)=f(x)\) 恒成立,且 \(f(x)\) 在 \(\left[0,\frac{\pi}{4}\right]\) 上有零点,则 \(\omega\) 的最小值为(\quad)
\choices
    {8}
    {6}
    {4}
    {3}
\end{problem}

\begin{problem}
在一定条件下,某大模型训练 \(N\) 个单位的数据量所需时间 \(T=k\log_2 N\)(单位:小时). 已知数据量从 \(10^6\) 增加到 \(1.024\times10^9\) 时增加 \(20\) 小时;当数据量从 \(1.024\times10^9\) 增加到 \(4.096\times10^9\) 时,训练时间增加(\quad)
\choices
    {2}
    {4}
    {20}
    {40}
\end{problem}

\begin{problem}
已知平面直角坐标系 \(xOy\) 中,\(|\overrightarrow{OA}|=|\overrightarrow{OB}|=\sqrt2\),\(|\overrightarrow{AB}|=2\),设 \(C(3,4)\),则 \(|2\overrightarrow{CA}+\overrightarrow{AB}|\) 的取值范围是(\quad)
\choices
    {\([6,14]\)}
    {\([6,12]\)}
    {\([8,14]\)}
    {\([8,12]\)}
\end{problem}

\section{填空题}

\begin{problem}
抛物线 \(y^2=2px\)(\(p>0\))的顶点到焦点的距离为 3,则 \(p=\)\fillinblank{}.
\end{problem}

\begin{problem}
已知 \((1-2x)^4=a_0-2a_1x+4a_2x^2-8a_3x^3+16a_4x^4\),则 \(a_0=\);\(a_1+a_2+a_3+a_4=\)\fillinblank{}.
\end{problem}

\begin{problem}
已知 \(\alpha\),\(\beta\in[0,2\pi]\),且 \(\sin(\alpha+\beta)=\sin(\alpha-\beta)\),\(\cos(\alpha+\beta)\ne\cos(\alpha-\beta)\),写出满足条件的一组 \(\alpha=\fillinblank{}\),\(\beta=\fillinblank{}\).
\end{problem}

\begin{problem}
某科技兴趣小组通过 3D 打印机打印的一个零件可以抽象为如图所示的多面体,其中 \(ABCDEF\) 是一个平行六边形,平面 \(ARF\perp\) 平面 \(ABC\),平面 \(TCD\perp\) 平面 \(ABC\),\(AB\perp BC\),\(AB\parallel RS\parallel EF\parallel CD\),\(AF\parallel ST\parallel BC\parallel ED\). 若 \(AB=BC=8\),\(AF=CD=4\),\(AR=RF=TC=TD=\dfrac52\),则该多面体的体积为(\quad)

\bitmapfigure[max width=0.48\linewidth]{img/2025/beijing/q14_fig1.png}

\choices
    {60}
    {90}
    {120}
    {150}
\end{problem}

\section{解答题}

\begin{problem}
关于定义域为 \(\R\) 的函数 \(f(x)\),以下说法正确的有(\quad)

\choices
    {存在在 \(\R\) 上单调递增的函数 \(f(x)\),使 \(f(x)+f(2x)=-x\) 恒成立}
    {存在在 \(\R\) 上单调递减的函数 \(f(x)\),使 \(f(x)+f(2x)=-x\) 恒成立}
    {存在且有无穷多个函数 \(f(x)\),满足 \(f(x)+f(-x)=\cos x\)}
    {存在且有无穷多个函数 \(f(x)\),满足 \(f(x)-f(-x)=\cos x\)}
\end{problem}

\begin{problem}
在 \(\triangle ABC\) 中,\(\cos A=-\dfrac13\),\(a\sin C=4\sqrt2\).

\begin{enumerate}
\item 求 \(c\);

\item 在以下三个条件中选择一个作为已知,使得 \(\triangle ABC\) 存在,求 \(BC\) 的高.
\end{enumerate}

\(\text{\circled{1}} \, a=6\),\(\text{\circled{2}} \, b\sin C= \frac{10\sqrt2}{3}\),\(\text{\circled{3}} \, S_{\triangle ABC}=10\sqrt2\)
\end{problem}

\begin{problem}
四棱锥 \(P-ABCD\),\(\triangle ACD\) 与 \(\triangle ABC\) 为等腰直角三角形,\(\angle ADC=90^\circ\),\(\angle BAC=90^\circ\),\(E\) 为 \(BC\) 中点.

\begin{enumerate}
\item \(F\) 为 \(PD\) 中点,\(G\) 为 \(PE\) 中点,证明 \(FG\parallel\) 平面 \(PAB\);

\item 若 \(PA\perp\) 平面 \(ABCD\),\(\ PA=AC\),求 \(AB\) 与平面 \(PCD\) 所成角的正弦值.
\end{enumerate}
\end{problem}

\begin{problem}
有一道选择题考查某知识点. 甲,乙两校各随机抽取 \(100\) 人,甲校有 \(80\) 人答对,乙校有 \(75\) 人答对,用频率估计概率.

\begin{enumerate}
    \item 从甲校随机抽取 1 人,求这个人做对该题目的概率;
    \item 从甲乙两校各随机抽取 1 人,设 \(X\) 为做对人数,求 \(P(X=1)\) 与 \(E(X)\);
    \item 若甲校同学掌握该知识点,则有 \(100\%\) 的概率做对该题目;乙校同学掌握该知识点,则有 \(85\%\) 的概率做对该题目;未掌握该知识点的同学都是从四个选项里面随机选择一个. 设甲校学生掌握该知识点的概率为 \(p_1\),乙校学生掌握该知识点的概率为 \(p_2\),比较 \(p_1\) 与 \(p_2\) 的大小.
\end{enumerate}
\end{problem}

\begin{problem}
已知椭圆 \(\dfrac{x^2}{a^2}+\dfrac{y^2}{b^2}=1\)(\(a>b>0\)),其离心率为 \(\dfrac{\sqrt{2}}{2}\),椭圆上一点到两焦点距离之和为 \(4\).

\begin{enumerate}
\item 求椭圆方程;
\item 设 \(O=(0,0)\),\(M(x_0,y_0)\)(\(x_0\ne0\))为椭圆上一点,直线 \(x_0x+2y_0y-4=0\) 与直线 \(y=2\),\(y=-2\) 交于 \(A\),\(B\). \(S_1\),\(S_2\) 分别为 \(\triangle OAM\) 与 \(\triangle OBM\) 的面积,比较 \(\dfrac{S_1}{S_2}\) 与 \(\dfrac{|OA|}{|OB|}\) 的大小.
\end{enumerate}
\end{problem}

\begin{problem}
函数 \(f:(-1,+\infty)\to\R\),\(f(0)=0\),\(f'(x)=\dfrac{\ln(1+x)}{1+x}\),\(l_1\) 为点 \(A(a,f(a))\)(\(a\ne0\))处切线.

\begin{enumerate}
    \item \(f'(x)\) 的最大值;
    \item 证明:当 \(-1<a<0\) 时,除点 \(A\) 外曲线 \(y=f(x)\) 在 \(l_1\) 上方;
    \item 若 \(a>0\) 时,过 \(A\) 且与 \(l_1\) 垂直的 \(l_2\) 与 \(l_1\) 分别交 \(x\) 轴于 \(x_1\),\(x_2\),求 \(\dfrac{2a-x_1-x_2}{x_2-x_1}\) 的范围.
\end{enumerate}
\end{problem}

\begin{problem}
设 \(M=\{(x_i,y_i)\mid x_i,y_i\in A\}\),\(A=\{1,2,3,4,5,6,7,8\}\). 从 \(M\) 中选 \(n\) 个有序数对构成一列 \((x_1,y_1)\),\(\dots\),\((x_n,y_n)\). 若相邻两项满足 \(|x_{i+1}-x_i|=3\text{\ 或\ }4\),\(|y_{i+1}-y_i|=4\text{\ 或\ }3\),\(i=1\),\(\dots\),\(n-1\),则该列称为 \(k\) 列.

\begin{enumerate}
    \item 若 \(k\) 列的第一项为 \((3,3)\),求第二项.
    \item 若 \(\tau\) 为 \(k\) 列,且 \(i\) 为奇数时 \(x_i\in\{1,2,7,8\}\),\(i\) 为偶数时 \(x_i\in\{3,4,5,6\}\),判断 \((3,2)\) 与 \((4,4)\) 能否同时在 \(\tau\) 中,并说明.
    \item 证明 \(M\) 中所有元素都不构成 \(k\) 列.
\end{enumerate}
\end{problem}

